% Сборник задач — глава 2
% Номер главы определяется именем этого файла.

\begin{exo}[От холодного льда до кипящей воды: оценки порядка величины]{chauffage-glace-eau-ebullition}
\seoready{true}
\keywords{калориметрия, фазовый переход, плавление, испарение, скрытая теплота, порядок величины}

\textbf{Числовые данные (при атмосферном давлении):}
\[
\begin{aligned}
c_{\text{ice}} &= 2{,}1\times10^3\,\text{J}\!\cdot\!\text{kg}^{-1}\!\cdot\!\text{K}^{-1},
& L_f &= 3{,}34\times10^5\,\text{J}\!\cdot\!\text{kg}^{-1},\\
c_{\text{water}} &= 4{,}18\times10^3\,\text{J}\!\cdot\!\text{kg}^{-1}\!\cdot\!\text{K}^{-1},
& L_v &= 2{,}26\times10^6\,\text{J}\!\cdot\!\text{kg}^{-1}.
\end{aligned}
\]
При атмосферном давлении постепенно нагревают один килограмм льда с начальной температурой $-100\,{}^\circ\text{C}$.

\begin{enumerate}
  \item Вычислить энергию, необходимую для нагревания льда от $-100\,{}^\circ\text{C}$ до $0\,{}^\circ\text{C}$.
  \item Вычислить энергию, необходимую для полного плавления льда при $0\,{}^\circ\text{C}$.
  \item Вычислить энергию, необходимую для нагревания жидкой воды от $0\,{}^\circ\text{C}$ до $100\,{}^\circ\text{C}$.
  \item Вычислить дополнительную энергию, необходимую для полного испарения воды при $100\,{}^\circ\text{C}$.
  \item Какой этап требует больше всего энергии? Сравнить энергию нагревания одного килограмма воды от $0\,{}^\circ\text{C}$ до $100\,{}^\circ\text{C}$ с потенциальной энергией массы $m$, падающей с высоты десяти метров. Какая масса должна упасть, чтобы выделилась такая же энергия?
  \item Рассмотреть обратный процесс. Масса $m_v$ водяного пара конденсируется на ветровом стекле, а полученная жидкая вода затем не охлаждается. Выразить теплоту $Q_{\text{rel}}$, выделившуюся при конденсации, и вычислить её для $m_v=1{,}0\times10^{-2}\,\text{kg}$. При температуре стекла принять $L_v=2{,}45\times10^6\,\text{J}\!\cdot\!\text{kg}^{-1}$.
\end{enumerate}

\begin{indication}
Для нагревания фазы без фазового перехода использовать $Q=mc\Delta T$. Для фазового перехода при постоянной температуре использовать $Q=mL$. Потенциальная энергия массы $m$ на высоте $h$ равна $E_p=mgh$; принять $g=9{,}81\,\text{m}\!\cdot\!\text{s}^{-2}$.
\end{indication}

\begin{solution}
\textbf{1.} Нагревание льда до температуры плавления требует
\[
Q_1=mc_{\text{ice}}\Delta T
=1{,}0\times2{,}1\times10^3\times100
=2{,}1\times10^5\,\text{J}.
\]
\textbf{2.} Во время плавления температура остаётся равной $0\,{}^\circ\text{C}$:
\[
Q_2=mL_f=1{,}0\times3{,}34\times10^5=3{,}34\times10^5\,\text{J}.
\]
\textbf{3.} Для нагревания жидкой воды
\[
Q_3=mc_{\text{water}}\Delta T
=1{,}0\times4{,}18\times10^3\times100
=4{,}18\times10^5\,\text{J}.
\]
\textbf{4.} Полное испарение дополнительно требует
\[
Q_4=mL_v=1{,}0\times2{,}26\times10^6=2{,}26\times10^6\,\text{J}.
\]
Один этот этап требует энергии порядка нескольких мегаджоулей.

\textbf{5.} Испарение — безусловно самый энергоёмкий этап:
\[
Q_4=2{,}26\times10^6>Q_3=4{,}18\times10^5>Q_2=3{,}34\times10^5>Q_1=2{,}1\times10^5\,\text{J}.
\]
Оно требует примерно $(2{,}26\times10^6)/(4{,}18\times10^5)\approx5{,}4$ раза больше энергии, чем нагревание жидкой воды от $0$ до $100\,{}^\circ\text{C}$.

Для падения с высоты $h=10\,\text{m}$ имеем $E_p=mgh$. Из $mgh=Q_3$ следует
\[
m=\frac{Q_3}{gh}=\frac{4{,}18\times10^5}{9{,}81\times10}
\approx4{,}26\times10^3\,\text{kg}.
\]
Чтобы выделить столько же энергии, сколько требуется для нагревания килограмма воды от $0$ до $100\,{}^\circ\text{C}$, с высоты десяти метров должна упасть масса около $4{,}3$ тонны. Это огромная величина, объясняющая существенную долю нагревания воды в бытовом энергопотреблении. Удельная теплоёмкость воды также очень велика по сравнению с обычными металлами: например, примерно в десять раз больше, чем у меди (см. следующие задачи).

\textbf{6.} Конденсация является процессом, обратным испарению: пар отдаёт стеклу скрытую теплоту
\[
Q_{\text{rel}}=m_vL_v.
\]
Для $m_v=1{,}0\times10^{-2}\,\text{kg}$,
\[
Q_{\text{rel}}=1{,}0\times10^{-2}\times2{,}45\times10^6
=2{,}45\times10^4\,\text{J}.
\]
Таким образом, даже небольшая масса конденсирующегося пара может выделить заметное количество теплоты, которое получают стекло и его окружение.
\end{solution}
\end{exo}

\begin{exo}[Эвапотранспирация большого дерева: природный кондиционер?]{odg-evapotranspiration-arbre}
\seoready{true}
\keywords{порядок величины, эвапотранспирация, дерево, скрытая теплота, холодильная мощность, кондиционер, COP}

В жаркий сухой день большое дерево при достаточном водоснабжении может испарить путём эвапотранспирации около $500\,\text{L}=5{,}0\times10^{-1}\,\text{m}^3$ воды. Поскольку транспирация происходит главным образом при освещении, будем считать, что она длится двенадцать часов. Крона дерева покрывает площадь земли $A_{\text{tree}}=50\,\text{m}^2$. Даны плотность воды и её удельная теплота парообразования при атмосферном давлении:
\[
\rho_{\text{water}}=1{,}0\times10^3\,\text{kg}\!\cdot\!\text{m}^{-3},
\qquad L_v=2{,}45\times10^6\,\text{J}\!\cdot\!\text{kg}^{-1}.
\]

\begin{enumerate}
  \item Вычислить массу воды, испарённой деревом за день, и энергию, необходимую для её испарения.
  \item Объяснить, почему это создаёт охлаждающий эффект.
  \item Вычислить среднюю охлаждающую мощность дерева на квадратный метр в течение двенадцати часов активной транспирации.
  \item Для сравнения рассмотреть кондиционер с электрической мощностью $P_{\mathrm{el}}=2{,}0\times10^3\,\text{W}$, холодильным коэффициентом $\mathrm{COP}=3{,}0$, охлаждающий помещение площадью $A_{\text{room}}=20\,\text{m}^2$. По определению $\mathrm{COP}=P_{\text{cool}}/P_{\mathrm{el}}$. Вычислить холодильную мощность и мощность на квадратный метр и сравнить их с результатами для дерева.
  \item Почему нельзя утверждать, что дерево и кондиционер создают совершенно одинаковое ощущение охлаждения, даже если их мощности на единицу площади имеют одинаковый порядок величины?
  \item Можно ли охлаждать жилище множеством комнатных растений? (Общие знания: ответ зависит от биологических процессов.)
\end{enumerate}

\begin{indication}
Использовать $m=\rho V$, $Q_{\mathrm{evap}}=mL_v$ и $P=Q/\tau$. Для кондиционера $P_{\text{cool}}=\mathrm{COP}\,P_{\mathrm{el}}$.
\end{indication}

\begin{solution}
\textbf{1.} Масса воды равна
\[
m=\rho_{\text{water}}V=1{,}0\times10^3\times5{,}0\times10^{-1}
=5{,}0\times10^2\,\text{kg}.
\]
Необходимая энергия составляет
\[
Q_{\mathrm{evap}}=mL_v=5{,}0\times10^2\times2{,}45\times10^6
=1{,}225\times10^9\,\text{J}.
\]
Её порядок величины — один гигаджоуль.

\textbf{2.} Испарение — эндотермический фазовый переход. Энергия забирается у листьев, почвы и окружающего воздуха, понижая температуру листьев и ближайшего окружения.

\textbf{3.} Двенадцать часов соответствуют $\tau=12\times3600=4{,}32\times10^4\,\text{s}$. Поэтому
\[
P_{\text{tree}}=\frac{Q_{\mathrm{evap}}}{\tau}
=\frac{1{,}225\times10^9}{4{,}32\times10^4}
\approx2{,}84\times10^4\,\text{W},
\]
а на единицу площади
\[
\frac{P_{\text{tree}}}{A_{\text{tree}}}
=\frac{2{,}84\times10^4}{50}
\approx5{,}7\times10^2\,\text{W}\!\cdot\!\text{m}^{-2}.
\]

\textbf{4.} Полезная холодильная мощность кондиционера равна
\[
P_{\text{cool}}=\mathrm{COP}\,P_{\mathrm{el}}
=3{,}0\times2{,}0\times10^3=6{,}0\times10^3\,\text{W}.
\]
На единицу площади помещения приходится
\[
\frac{P_{\text{cool}}}{A_{\text{room}}}
=\frac{6{,}0\times10^3}{20}
=3{,}0\times10^2\,\text{W}\!\cdot\!\text{m}^{-2}.
\]
В этой модели мощность эвапотранспирации дерева на единицу площади почти вдвое больше, чем у обычного кондиционера.

\textbf{5.} Сравнение касается только потоков энергии. Кондиционер охлаждает замкнутый контролируемый объём, тогда как пар от дерева рассеивается в атмосфере, а ветер приносит тёплый воздух. Ощущаемый эффект зависит от вентиляции, влажности, температуры воздуха и расстояния до дерева. Транспирация также зависит от солнечного света, ветра, влажности, вида дерева и доступности воды и может резко уменьшиться при засухе. Кроме того, дерево создаёт тень и уменьшает солнечное излучение, получаемое землёй и людьми. Расчёт сравнивает порядки величины, но не гарантирует одинакового теплового комфорта.

\textbf{6.} На практике охлаждающий эффект комнатных растений очень мал. У большинства растений свет способствует открытию устьиц — мелких пор листа, через которые выходит водяной пар. Освещение в помещении обычно слабое, поэтому устьица открываются меньше и транспирация снижается. В плохо проветриваемой комнате рост влажности также постепенно замедляет испарение.

Множество растений может создать небольшое местное испарительное охлаждение, но не заменит кондиционер. Для длительного удаления энергии влажный воздух пришлось бы выводить наружу. Интенсивное искусственное освещение могло бы усилить транспирацию, однако его электрическая энергия почти полностью превратилась бы в тепло в комнате. Приспособленные к засухе CAM-растения, включая кактусы, являются исключением: их устьица открываются главным образом ночью.
\end{solution}
\end{exo}

\begin{exo}[Смешивание двух масс воды]{calorimetrie-melange-eau}
\seoready{true}
\keywords{калориметрия, смешивание, теплоёмкость, удельная теплоёмкость, Джозеф Блэк, температура равновесия}

$m_1=0{,}200\,\text{kg}$ воды при $T_1=80\,^\circ\text{C}$ наливают в сосуд, уже содержащий $m_2=0{,}300\,\text{kg}$ воды при $T_2=15\,^\circ\text{C}$. Сосуд является идеальным калориметром: потерями теплоты в окружающую среду и теплоёмкостью самого сосуда пренебрегают. Напомним, что $Q=mc\Delta T$, где $c=4{,}18\times10^3\,\text{J}\!\cdot\!\text{kg}^{-1}\!\cdot\!\text{K}^{-1}$ — удельная теплоёмкость жидкой воды.

\begin{enumerate}
  \item Обосновать, что теплота, отданная горячей водой, полностью получена холодной водой. Записать уравнение сохранения с $m_1$, $m_2$, $c$, $T_1$, $T_2$ и температурой равновесия $T_{eq}$.
  \item Вывести выражение для $T_{eq}$ и вычислить численное значение.
  \item Вычислить теплоту $Q$, отданную горячей водой при смешивании.
  \item Позволяет ли смешивание одного и того же вещества определить $c$? Обосновать.
\end{enumerate}
\begin{indication}
Калориметр изолирован и жёсток, поэтому полная внутренняя энергия $U_1+U_2$ сохраняется: теплота, отданная одной частью, получена другой с противоположным знаком.
\end{indication}
\begin{solution}
\textbf{1.} Сохранение энергии даёт
\[
m_1c(T_{eq}-T_1)+m_2c(T_{eq}-T_2)=0.
\]
\textbf{2.} Общий множитель $c$ сокращается:
\[
T_{eq}=\frac{m_1T_1+m_2T_2}{m_1+m_2}
=\frac{0{,}200\times80+0{,}300\times15}{0{,}200+0{,}300}=41\,^\circ\text{C}.
\]
\textbf{3.} Горячая вода отдаёт
\[
Q=m_1c(T_1-T_{eq})=0{,}200\times4{,}18\times10^3\times(80-41)
\approx3{,}26\times10^4\,\text{J}.
\]
Холодная вода получает столько же:
\[
m_2c(T_{eq}-T_2)=0{,}300\times4{,}18\times10^3\times26
\approx3{,}26\times10^4\,\text{J}.
\]
\textbf{4.} Нет. Для двух масс одного вещества коэффициент $c$ умножает оба члена баланса и сокращается. Температура равновесия не зависит от $c$ и является средним начальных температур, взвешенным по массам. Для измерения удельной теплоёмкости методом смешивания нужно вещество с известной теплоёмкостью, как в следующей задаче.
\end{solution}
\end{exo}

\begin{exo}[Определение удельной теплоёмкости металла методом смешивания]{calorimetrie-methode-melanges}
\seoready{true}
\keywords{калориметрия, метод смешивания, удельная теплоёмкость, калориметр, водяной эквивалент}

Как это делал Джозеф Блэк уже в XVIII веке, требуется измерить удельную теплоёмкость $c_m$ неизвестного металла \emph{методом смешивания}. Образец массой $m=0{,}150\,\text{kg}$ нагревают до $T_m=95\,^\circ\text{C}$ и быстро погружают в калориметр с $m_e=0{,}250\,\text{kg}$ воды, для которой $c_e=4{,}18\times10^3\,\text{J}\!\cdot\!\text{kg}^{-1}\!\cdot\!\text{K}^{-1}$ и $T_e=18\,^\circ\text{C}$. Измеренная температура равновесия равна $T_{eq}=22\,^\circ\text{C}$. Сначала теплоёмкостью сосуда, мешалки и термометра пренебрегают.

\begin{enumerate}
  \item Записать сохранение энергии для изолированной системы «металл–вода», оставив $c_m$ единственной неизвестной.
  \item Вывести выражение для $c_m$ и вычислить его. Какому распространённому металлу соответствует результат? Справочные значения в $\text{J}\!\cdot\!\text{kg}^{-1}\!\cdot\!\text{K}^{-1}$: алюминий $9{,}0\times10^2$, железо $4{,}5\times10^2$, медь $3{,}9\times10^2$, свинец $1{,}3\times10^2$.
  \item Сосуд поглощает часть теплоты. Этот эффект описывается \emph{водяным эквивалентом} $\mu=1{,}5\times10^{-2}\,\text{kg}$: калориметр ведёт себя как дополнительная масса воды $\mu$. Повторно вычислить $c_m$ и объяснить направление поправки.
  \item Почему образец необходимо погружать быстро, а калориметр хорошо изолировать от окружающего воздуха?
\end{enumerate}
\begin{indication}
Металл отдаёт $Q_m=mc_m(T_m-T_{eq})$, полностью получаемое водой и, в вопросе 3, калориметром: $Q_m=(m_e+\mu)c_e(T_{eq}-T_e)$.
\end{indication}
\begin{solution}
\textbf{1.} Для изолированной системы
\[
mc_m(T_m-T_{eq})=m_ec_e(T_{eq}-T_e).
\]
\textbf{2.} Следовательно,
\[
c_m=\frac{m_ec_e(T_{eq}-T_e)}{m(T_m-T_{eq})}
=\frac{0{,}250\times4{,}18\times10^3\times(22-18)}{0{,}150\times(95-22)}
\approx3{,}82\times10^2\,\text{J}\!\cdot\!\text{kg}^{-1}\!\cdot\!\text{K}^{-1}.
\]
Это значение очень близко к значению для меди.

\textbf{3.} С учётом водяного эквивалента
\[
c_m=\frac{(m_e+\mu)c_e(T_{eq}-T_e)}{m(T_m-T_{eq})}
=\frac{(0{,}250+0{,}015)\times4{,}18\times10^3\times4}{0{,}150\times73}
\approx4{,}05\times10^2\,\text{J}\!\cdot\!\text{kg}^{-1}\!\cdot\!\text{K}^{-1}.
\]
Поправка увеличивает $c_m$: без учёта сосуда занижались реально отданная металлом теплота и его удельная теплоёмкость.

\textbf{4.} Метод предполагает изолированную систему в течение всего измерения. Быстрое погружение ограничивает теплообмен горячего образца с воздухом до попадания в воду, а хорошая изоляция уменьшает потери при установлении равновесия. Любой неучтённый поток теплоты искажает баланс и найденное $c_m$. Именно такую экспериментальную тщательность проявляли Блэк, а позднее Лавуазье и Лаплас в своём ледяном калориметре.
\end{solution}
\end{exo}

\begin{exo}[Пищевая калорийность и метаболическая мощность]{calorimetrie-calories-alimentaires}
\seoready{true}
\keywords{калория, килокалория, пищевая Калория, механический эквивалент, метаболическая мощность, единицы}

На пищевой этикетке указано, что взрослый человек в среднем потребляет около $2000\,\text{kcal}$ в сутки. Килокалория — внесистемная единица, всё ещё используемая в диетологии; $1\,\text{kcal}=4{,}18\times10^3\,\text{J}$.
\begin{enumerate}
  \item Перевести суточное потребление энергии в джоули.
  \item Определить среднюю метаболическую мощность в ваттах, считая расход энергии равномерным в течение 24 часов.
  \item На энергетическом батончике указано «$210\,\text{kcal}$». Какую массу воды с начальной температурой $20\,^\circ\text{C}$ можно было бы довести до кипения при $100\,^\circ\text{C}$ этой энергией при полном превращении её в теплоту? Принять $c_{\text{water}}=4{,}18\times10^3\,\text{J}\!\cdot\!\text{kg}^{-1}\!\cdot\!\text{K}^{-1}$.
  \item В каких формах организм человека фактически использует эту энергию? Дать качественный ответ.
\end{enumerate}
\begin{indication}
В вопросе 2 использовать $\mathcal P=E/\Delta t$. В вопросе 3 сначала перевести энергию в джоули, затем выразить $m$ из $Q=mc\Delta T$.
\end{indication}
\begin{solution}
\textbf{1.} В единицах СИ
\[
E=2000\times4{,}18\times10^3=8{,}36\times10^6\,\text{J}.
\]
\textbf{2.} Средняя мощность равна
\[
\mathcal P=\frac{8{,}36\times10^6}{8{,}64\times10^4}\approx97\,\text{W}.
\]
За сутки человек в среднем использует мощность, сравнимую с мощностью лампы накаливания. Этот часто неожиданный порядок величины показывает пользу перевода пищевых калорий в привычные физические единицы.

\textbf{3.} Энергия батончика равна
\[
Q=210\times4{,}18\times10^3\approx8{,}78\times10^5\,\text{J}.
\]
При $\Delta T=80\,\text{K}$,
\[
m=\frac{Q}{c\Delta T}
=\frac{8{,}78\times10^5}{4{,}18\times10^3\times80}
\approx2{,}63\,\text{kg}.
\]
Энергии одного батончика по порядку величины хватило бы, чтобы довести до кипения около $2{,}5\,\text{kg}$ водопроводной воды.

\textbf{4.} Организм не «расходует» энергию в смысле её исчезновения, а преобразует химическую энергию питательных веществ. Часть идёт на работу мышц и органов, поддержание электрических градиентов клеток, перенос молекул и химический синтез. Часть может запасаться как гликоген или жир. В конечном итоге большая часть использованной энергии рассеивается в виде теплоты, поддерживая температуру тела, а затем передаётся окружающей среде. Небольшая доля покидает организм как химическая энергия в отходах.
\end{solution}
\end{exo}
